\section{Modelo de Cálculo}

\subsection{Geometria e massas fixas}

A cabeça de guerra foi mantida com \(165~\text{kg}\). Seu comprimento foi estimado por volume equivalente:

\begin{equation}
    L_{CDG} = \frac{M_{CDG}}{\rho_{CDG}S_{ref}}
\end{equation}

Diferentemente da versão anterior, o comprimento total do míssil foi fixado em \(5{,}8~\text{m}\). Assim, depois de dimensionar booster, sustainer e tubeira, a seção de controle fecha o comprimento restante:

\begin{equation}
    L_{cont} = L - (L_{CDG}+L_B+L_S+L_{tub})
\end{equation}

A massa da seção de controle fecha a massa total:

\begin{equation}
    M_{cont}=M_T-(M_{pB}+M_{vasoB})-(M_{pS}+M_{vasoS})-M_{CDG}
\end{equation}

\subsection{Booster}

O booster acelera o míssil de \(V_i=0\) até \(V_f=0{,}9a\), com tempo de queima imposto de \(4~\text{s}\):

\begin{equation}
    F_B = M_T\frac{V_f - V_i}{T_{qB}}
\end{equation}

A massa de propelente é dada por:

\begin{equation}
    M_{pB}=\frac{F_BT_{qB}}{I_{spB}g_0}
\end{equation}

O código atualizado usa grão em estrela. A espessura queimada é:

\begin{equation}
    R_B = T_{qB}V_{qB}
\end{equation}

O comprimento do booster vem do volume de propelente dividido pela área efetiva da geometria em estrela:

\begin{equation}
    L_B=\frac{V_{pB}}{\pi R_B(\phi_{EB}-R_B)+N_{star}A_{star}}
\end{equation}

Também foi incluída a massa do vaso metálico do booster, estimada com espessura de \(5~\text{mm}\).

\subsection{Sustainer}

No cruzeiro, o sustainer compensa o arrasto:

\begin{equation}
    D = \frac{1}{2}\rho V^2S_{ref}C_D
\end{equation}

O \(C_D\) foi obtido por interpolação da tabela de domo ogival usada no código-base:

\begin{table}[H]
\centering
\caption{Tabela de arrasto adotada.}
\begin{tabular}{cc}
\toprule
Mach & \(C_D\) \\
\midrule
0{,}0 & 0{,}33 \\
0{,}7 & 0{,}33 \\
0{,}9 & 0{,}42 \\
1{,}1 & 0{,}67 \\
1{,}2 & 0{,}68 \\
1{,}5 & 0{,}70 \\
1{,}8 & 0{,}70 \\
2{,}0 & 0{,}67 \\
2{,}5 & 0{,}62 \\
3{,}0 & 0{,}61 \\
3{,}5 & 0{,}61 \\
\bottomrule
\end{tabular}
\end{table}

Para o sustainer, o código atualizado dimensiona a área de queima para produzir empuxo igual ao arrasto:

\begin{equation}
    S_S = \frac{D}{I_{spS}g_0\rho_SV_{qS}}
\end{equation}

O tempo de queima vem do alcance requerido:

\begin{equation}
    T_{qS}=\frac{70000}{V}
\end{equation}

E a massa de propelente do sustainer passa a ser calculada pela geometria e pelo tempo de queima, não mais por diferença de massa:

\begin{equation}
    M_{pS}=S_SL_S\rho_S, \qquad L_S=T_{qS}V_{qS}
\end{equation}

Também foi incluída a massa do vaso metálico do sustainer.
